% =============================================================================
% 5 MODELLBILDUNG UND PLAUSIBILISIERUNG
% Weitgehend ausgearbeitete Arbeitsfassung vom 07.08.2026
% =============================================================================
%
% Diese Fassung bildet den tatsächlichen Entwicklungsweg ab.
% Noch offene Angaben sind ausschließlich als LaTeX-Kommentare gekennzeichnet.
% Die auskommentierten Abbildungen verursachen keine Fehler, solange die
% zugehörigen Bilddateien noch nicht im Projektordner liegen.
%
% Morgen vor der endgültigen Übernahme noch abzugleichen:
% - genaue Implementierung des OSG im finalen Quellenmodell,
% - Bezugspunkt und Endwert der Quellenspannung,
% - Verwendung des phasenbezogenen oder gemittelten R_MCC im Gesamtmodell,
% - Solver, Schrittweiten und Simulationszeiten der endgültigen Modelle,
% - Export des zusammengeführten Gesamtmodells und seiner Ergebnisplots.
%
% =============================================================================

\section{Modellbildung und Plausibilisierung}
\label{chap:modellbildung}

Zur Untersuchung der in Kapitel~\ref{chap:analyse_pruefstand}
beschriebenen Stromabweichungen wird der betrachtete Lastzweig in
MATLAB/Simulink abgebildet. Das Ziel der Modellbildung besteht nicht darin,
sämtliche konstruktiven Einzelheiten des Erwärmungsprüfstands nachzubilden.
Vielmehr werden die elektrischen und thermischen Zusammenhänge erfasst, die
den Strom eines MCC-Abgangs und dessen zeitliche Änderung bestimmen. Das
zusammengeführte Gesamtmodell bildet anschließend die Referenz für die
Untersuchung der in Kapitel~\ref{chap:konzeptentwicklung} entwickelten
Stellkonzepte.

Die Modellbildung erfolgte nicht in einem einzigen Schritt. Zunächst wurde
ein vollständiger elektrischer Signal- und Energiepfad aufgebaut und mit
vorläufigen Parametern erweitert. Da mehrere noch nicht überprüfte
Teilfunktionen gleichzeitig zusammenwirkten, konnten die beobachteten
Abweichungen und Schwingungen keiner eindeutigen Ursache zugeordnet werden.
Das Ausgangsmodell wurde deshalb in die Teilbereiche Spannungsquelle,
elektrischer MCC-Pfad und thermisches Verhalten des V2A-Widerstands zerlegt.
Erst nach der getrennten Untersuchung und Plausibilisierung dieser
Teilmodelle werden sie zu einem Referenzmodell des bestehenden Lastzweigs
zusammengeführt.


% -----------------------------------------------------------------------------
% 5.1 Ausgangsmodell und methodische Zerlegung
% -----------------------------------------------------------------------------

\subsection{Ausgangsmodell und methodische Zerlegung}
\label{sec:ausgangsmodell_zerlegung}

Der erste Modellstand entstand als funktionaler Aufbau ohne zuvor
vollständig festgelegte Modellhierarchie. Zunächst wurde die
Spannungsquelle erstellt. An diese wurde ein repräsentativer Lastzweig des
\SI{400}{\ampere}-Moduls angeschlossen. Der Widerstand der V2A-Schienen wurde
dabei nicht frei gewählt, sondern aus deren Geometrie und dem spezifischen
elektrischen Widerstand berechnet. Demgegenüber wurde der Innenwiderstand des
MCC-Moduls in dieser frühen Phase lediglich als pauschaler Widerstand
angenommen. Für diesen Parameter lagen zu diesem Zeitpunkt noch keine
Messwerte vor.

Der wirksame V2A-Widerstand wurde zunächst mithilfe eines Sliders manuell
verändert. Dadurch ließ sich prüfen, ob die Verschaltung grundsätzlich auf
eine Widerstandsänderung mit der erwarteten Stromänderung reagiert. Im
nächsten Schritt ersetzte eine exponentielle Testfunktion die manuelle
Verstellung. Diese Funktion stellte noch kein thermisches Modell dar. Sie
diente ausschließlich dazu, eine zeitlich veränderliche Störgröße zu
erzeugen und die Reaktion des elektrischen Modells zu untersuchen.

Anschließend wurde der Lastzweig um den parallelen Pfad aus Trafo~2 und
Stellwiderstand erweitert. Auf dieser Grundlage wurden erste geschlossene
Regelkreise untersucht. Die dabei auftretenden Schwingungen konnten jedoch
nicht eindeutig dem Regler zugeordnet werden. Neben der Reglerparametrierung
wirkten gleichzeitig die angenommene Modulimpedanz, die künstliche
Widerstandsänderung, die Quellen- und Messwertaufbereitung sowie der
zusätzliche Strompfad. Eine Bewertung der Regelgüte wäre auf dieser Grundlage
nicht belastbar gewesen.

Diese Erkenntnis führte zur Zerlegung des Modells. Die Zerlegung ist damit
kein nachträglich gewähltes Darstellungsprinzip, sondern das Ergebnis der
ersten Simulationsversuche. Tabelle~\ref{tab:modellstufen_kap5} fasst die
daraus entstandenen Modellstufen zusammen.

\begin{table}[H]
    \centering
    \caption{Entwicklungsstufen der Modellbildung}
    \label{tab:modellstufen_kap5}
    \begin{tabularx}{\textwidth}{@{}p{0.20\textwidth}p{0.28\textwidth}X@{}}
        \toprule
        Stufe & Inhalt & Zweck \\
        \midrule
        Ausgangsmodell &
        Quelle, angenommener \(R_\mathrm{MCC}\), geometrisch bestimmter
        V2A-Widerstand und repräsentativer \SI{400}{\ampere}-Abgang &
        Prüfung der grundsätzlichen Verschaltung und der Signalrichtungen \\

        Erweiterung &
        Slider, exponentielle Teststörung, paralleler Stellpfad und Regler &
        Untersuchung der Stellwirkung und Erkennen der nicht eindeutig
        zuordenbaren Abweichungen \\

        Quellenmodell &
        sinusförmige Momentanspannung, lineare Effektivwertrampe,
        RMS-Auswertung und OSG &
        eindeutige Festlegung der elektrischen Anregung und der verwendeten
        Effektivwertgröße \\

        Elektrisches Teilmodell &
        Messung und Berechnung des Ersatzwiderstands
        \(R_\mathrm{MCC}\) &
        Ersatz der vorläufigen Widerstandsannahmen durch messgestützte
        Modellparameter \\

        Thermisches Teilmodell &
        Kreis~1, Kreis~2 und Kreis~3 &
        schrittweise Prüfung des thermischen Grundmodells, seiner
        Parametrierung und der elektrothermischen Rückkopplung \\

        Gesamtmodell &
        Zusammenführung der linearen Quelle, des gemessenen
        \(R_\mathrm{MCC}\) und des elektrothermischen Modells &
        Untersuchung des Stromdrifts unter den festgelegten Randbedingungen \\
        \bottomrule
    \end{tabularx}
\end{table}

% BILD 5.1 -- erst aktivieren, sobald der Export des ersten Gesamtmodells
% vorliegt.
% \begin{figure}[H]
%     \centering
%     \includegraphics[width=\textwidth]
%     {03_Ressourcen/pic/Ausgangsmodell_freier_Aufbau.pdf}
%     \caption{Erster funktionaler Modellstand vor der Zerlegung in
%     überprüfbare Teilmodelle (eigene Darstellung in MATLAB/Simulink)}
%     \label{fig:kap5_ausgangsmodell}
% \end{figure}

Die nachfolgende Darstellung folgt dieser tatsächlichen Entwicklung. Die
frühen Annahmen werden dabei nicht als Parameter des endgültigen Modells
übernommen. Sie werden nur dort beschrieben, wo sie die Entscheidung zur
Zerlegung und Weiterentwicklung des Modells erklären.


% -----------------------------------------------------------------------------
% 5.2 Modellierung der Spannungsquelle
% -----------------------------------------------------------------------------

\subsection{Modellierung der Spannungsquelle und Effektivwertaufbereitung}
\label{sec:spannungsquelle_osg}

Der Säulenstelltransformator des realen Prüfstands verändert die
Prüfspannung mechanisch. Die sekundärseitige Spannung erreicht ihren Endwert
daher nicht sprunghaft, sondern steigt während der Stellbewegung an. Für das
Quellenmodell wird dieser Vorgang zunächst durch einen linearen Anstieg des
Phaseneffektivwerts angenähert. Die Spannung ist dabei als Spannung eines
Außenleiters gegen den gemeinsamen Sternpunkt des Prüfaufbaus definiert. Da
die drei Außenleiter im späteren Prüfstandsbetrieb einzeln ausgewertet und
nachgeführt werden, wird die Signalaufbereitung für jede Phase separat
angewendet. Eine Symmetrie der drei Phasenspannungen wird für die
Effektivwertbestimmung somit nicht vorausgesetzt.

Für eine betrachtete Phase lautet die sinusförmige Momentanspannung

\begin{equation}
    u_{\alpha}(t)
    =
    \sqrt{2}\,U_\mathrm{eff}(t)
    \sin\!\left(2\pi f t+\varphi_0\right),
    \qquad f=\SI{50}{\hertz}.
    \label{eq:kap5_momentanspannung}
\end{equation}

Für den untersuchten Quellenvergleich wird der Effektivwert durch

\begin{equation}
    U_\mathrm{eff}(t)
    =
    \begin{cases}
        U_\mathrm{eff,max}\dfrac{t}{t_\mathrm{r}},
        & 0\leq t<t_\mathrm{r}, \\[3mm]
        U_\mathrm{eff,max},
        & t\geq t_\mathrm{r}
    \end{cases}
    \label{eq:kap5_spannungsrampe}
\end{equation}

beschrieben. In den bisher durchgeführten Quellenvergleichen wurden
\(U_\mathrm{eff,max}=\SI{4.0}{\volt}\) und
\(t_\mathrm{r}=\SI{5}{\second}\) verwendet. Diese Werte kennzeichnen den
untersuchten Simulationsfall. Insbesondere der endgültige Spannungsendwert des
gekoppelten Modells ist später an den Bezugspunkt der Spannung und die dann
berücksichtigten Serienwiderstände anzupassen.

Eine klassische Effektivwertbildung bestimmt den quadratischen Mittelwert
über ein Zeitfenster der Länge \(T=1/f\):

\begin{equation}
    U_\mathrm{RMS}(t)
    =
    \sqrt{
        \frac{1}{T}
        \int_{t-T}^{t}u_{\alpha}^{2}(\lambda)\,\mathrm{d}\lambda
    }.
    \label{eq:kap5_rms_fenster}
\end{equation}

Für eine stationäre sinusförmige Spannung liefert diese Beziehung den
korrekten Effektivwert. Während der Rampe enthält das Messfenster jedoch
Momentanwerte mit unterschiedlichen Amplituden. Der berechnete RMS-Wert folgt
der aktuellen Amplitude deshalb zeitlich verzögert. Dieses Verhalten ist kein
Fehler der Effektivwertdefinition, erschwert aber die eindeutige Auswertung
des vorgegebenen linearen Hochlaufs.

Aus diesem Grund wurde zusätzlich eine orthogonale Signalgenerierung
verwendet. Zu \(u_{\alpha}\) wird eine um \(90^\circ\) verschobene Komponente
\(u_{\beta}\) gebildet. Die Bezeichnungen \(\alpha\) und \(\beta\) beziehen
sich hier auf das orthogonale Signalpaar einer einzelnen Phase. Es handelt
sich nicht um eine Clarke-Transformation, bei der ein dreiphasiges System
gemeinsam ausgewertet wird. Im idealen Fall gilt

\begin{align}
    u_{\alpha}(t)
    &=
    \sqrt{2}\,U_\mathrm{eff}(t)\sin(\omega t),
    \\
    u_{\beta}(t)
    &=
    \sqrt{2}\,U_\mathrm{eff}(t)\cos(\omega t).
    \label{eq:kap5_osg_komponenten}
\end{align}

Der Betrag des daraus gebildeten Effektivwertzeigers ergibt sich zu

\begin{equation}
    U_\mathrm{eff,Zeiger}(t)
    =
    \frac{1}{\sqrt{2}}
    \sqrt{u_{\alpha}^{2}(t)+u_{\beta}^{2}(t)}.
    \label{eq:kap5_osg_betrag}
\end{equation}

Bei idealer Orthogonalität entspricht dieser Wert unmittelbar der
vorgegebenen Effektivwertamplitude. Im Vergleich mit der klassischen
RMS-Auswertung folgt der Effektivwertzeiger der linearen Rampe ohne die durch
das gleitende Messfenster verursachte Nachwirkung. Für das elektrische
Lastmodell wird deshalb der Effektivwertzeiger verwendet. Die
Momentanspannung bleibt parallel verfügbar, damit Spannungsform, Frequenz und
Scheitelwert getrennt kontrolliert werden können. Derselbe Aufbau kann für
L1, L2 und L3 verwendet werden, wobei für jede Phase ein eigener
Effektivwertzeiger ausgegeben wird.

% OFFEN:
% Im endgültigen Text anhand des Simulink-Modells ergänzen, ob u_beta direkt
% in der MATLAB-Funktion oder in einem eigenständigen OSG-Subsystem erzeugt
% wurde. Keine konkrete OSG-Struktur behaupten, bevor das finale Modellbild
% und der zugehörige Code abgeglichen wurden.

% BILD 5.2 -- Quellenmodell mit Momentanspannung und OSG.
% \begin{figure}[H]
%     \centering
%     \includegraphics[width=\textwidth]
%     {03_Ressourcen/pic/Spannungsquelle_lineare_Rampe_OSG.pdf}
%     \caption{Modell der sinusförmigen Spannungsquelle mit linearer
%     Effektivwertrampe und orthogonaler Signalgenerierung
%     (eigene Darstellung in MATLAB/Simulink)}
%     \label{fig:kap5_spannungsquelle_osg}
% \end{figure}

% BILD 5.3 -- Vergleich RMS und Zeiger.
% \begin{figure}[H]
%     \centering
%     \includegraphics[width=\textwidth]
%     {03_Ressourcen/plots/02_Vergleich_RMS_vs_Zeiger.pdf}
%     \caption{Vergleich der klassischen RMS-Auswertung mit dem
%     Effektivwertzeiger während des linearen Spannungshochlaufs}
%     \label{fig:kap5_vergleich_rms_zeiger}
% \end{figure}

% BILD 5.4 -- Vergleich Zeiger und ideale Rampe.
% \begin{figure}[H]
%     \centering
%     \includegraphics[width=\textwidth]
%     {03_Ressourcen/plots/03_Vergleich_Zeiger_vs_Ideal.pdf}
%     \caption{Vergleich des Effektivwertzeigers mit der ideal vorgegebenen
%     linearen Effektivwertrampe}
%     \label{fig:kap5_vergleich_zeiger_ideal}
% \end{figure}

Die lineare Rampe ist von dem später im thermischen Teilmodell verwendeten
Spannungssprung zu unterscheiden. Der Sprung dient dort als definierte
Prüfanregung zur Untersuchung des thermischen Einschwingvorgangs. Er bildet
nicht das reale Verfahren des Säulenstelltransformators ab und wird daher im
zusammengeführten Gesamtmodell wieder durch das Quellenmodell nach
Gleichung~\eqref{eq:kap5_spannungsrampe} ersetzt.


% -----------------------------------------------------------------------------
% 5.3 Messung und Modellierung von R_MCC
% -----------------------------------------------------------------------------

\subsection{Messung und Modellierung des MCC-Ersatzwiderstands}
\label{sec:rmcc_messung_modellierung}

Im ersten Modellstand wurde der Widerstand der MCC-Abgangseinheit mangels
Messwerten pauschal angenommen. Bei den sehr niedrigen Prüfspannungen besitzt
dieser Anteil jedoch einen wesentlichen Einfluss auf den Laststrom. Die
Ermittlung des MCC-Ersatzwiderstands erfolgte in drei zeitlich getrennten
Messständen. Diese Messstände besitzen nicht dieselbe Aussagekraft und werden
deshalb im Folgenden nicht zu einer gemeinsamen Messreihe zusammengefasst.

\noindent\textbf{Erste Messung: orientierende Untersuchung des Messaufbaus.}
Die erste Messung wurde ohne Anwesenheit des Verfassers durchgeführt. Die
übermittelte Dokumentation nennt für ein 1er-Modul einen Prüfstrom von
\SI{100.2}{\ampere} und für ein 3er-Modul einen Prüfstrom von
\SI{300}{\ampere}. Da der genaue Aufbau, die Lage sämtlicher
Spannungsabgriffpunkte und die Führung der Messleitungen nicht vollständig
nachvollzogen werden konnten, werden die dabei ermittelten Spannungswerte
nicht zur Parametrierung des Modells verwendet.

Aus der ersten Messung ließ sich dennoch das Messprinzip ableiten. Der
Prüfstrom wurde über die Leistungskontakte des
ausgebauten MCC-Einschubs geführt. Zusätzliche Kupferstücke verbanden die
    Kontaktstellen so, dass die zu untersuchenden Phasenpfade einzeln oder in
Reihe bestromt werden konnten. Dadurch durchfloss der Prüfstrom nicht nur die
eingebauten Schaltgeräte und internen Leiter, sondern auch die innerhalb der
gewählten Spannungs\-abgriffpunkte liegenden Steck- und
Übergangs\-kontakte. Der
zu bestimmende Ersatzwiderstand beschreibt somit den eingebauten
MCC-Strompfad und nicht lediglich die Summe einzelner
Datenblatt\-widerstände.

\noindent\textbf{Zweite Messung: fotografische Erfassung des Einschubs.}
Auch der zweite Messstand war für eine quantitative Widerstandsbestimmung
nicht hinreichend reproduzierbar dokumentiert. Seine Messwerte werden daher
ebenfalls nicht in das Modell übernommen. Die dabei entstandenen Aufnahmen
sind jedoch für die konstruktive Beschreibung des MCC-Einschubs geeignet.
Sie zeigen die rückseitigen Leistungskontakte, die internen Kupferleiter, das
Schutz- und Schaltgerät sowie die SIMOCODE-Komponenten. Der Hauptstrompfad
verläuft damit vereinfacht von den einspeiseseitigen Leistungskontakten über
die internen Leiter und die Schutz- beziehungsweise Schaltgeräte, durch die
Stromerfassung und anschließend zu den abgangsseitigen Leistungskontakten.
Die genaue Reihenfolge ist abhängig von der jeweiligen Modulbestückung.

% BILD RMCC-A -- Bauteilübersicht aus der zweiten Messung.
% Geeignete Quellen: Mess 2.pdf, Seite 3 (Gesamtansicht),
% Seite 4 (SIRIUS-Schutzschalter) und Seite 5 (SIMOCODE PRO).
% \begin{figure}[H]
%     \centering
%     \begin{subfigure}[t]{0.32\textwidth}
%         \centering
%         \includegraphics[width=\textwidth]
%         {03_Ressourcen/pic/MCC_Einschub_Innenansicht.jpg}
%         \caption{Innenansicht}
%     \end{subfigure}
%     \hfill
%     \begin{subfigure}[t]{0.32\textwidth}
%         \centering
%         \includegraphics[width=\textwidth]
%         {03_Ressourcen/pic/MCC_Schutzschalter.jpg}
%         \caption{Schutzschalter}
%     \end{subfigure}
%     \hfill
%     \begin{subfigure}[t]{0.32\textwidth}
%         \centering
%         \includegraphics[width=\textwidth]
%         {03_Ressourcen/pic/MCC_SIMOCODE.jpg}
%         \caption{SIMOCODE-Komponente}
%     \end{subfigure}
%     \caption{Ausgewählte Bauteile und Leitungsführung innerhalb eines
%     MCC-Einschubs}
%     \label{fig:kap5_mcc_bauteile}
% \end{figure}

\noindent\textbf{Dritte Messung: Parameterbestimmung für das Modell.}
Erst die dritte Messung wurde mit einem einheitlichen Messablauf durchgeführt
und sowohl fotografisch als auch in einem strukturierten Messprotokoll
dokumentiert. Die Fotodokumentation hält für die einzelnen Messwerte den
Prüfstrom, die Lage der Spannungsspitzen, die Kupferverbindungen und die
jeweilige Anzeige des Millivoltmeters fest. Das zugehörige Tabellenprotokoll
ordnet diese Einzelaufnahmen den Messpunkten zu und berechnet daraus
Widerstände und Verlustleistungen. Erst die gemeinsame Betrachtung beider
Unterlagen macht den Messwert räumlich und rechnerisch nachvollziehbar.
Das vollständige Messprotokoll ist deshalb Bestandteil der
Versuchsdokumentation. Im Haupttext werden daraus die für die Modellbildung
maßgebenden Gesamt-, Phasen- und Teilspannungsfälle zusammengefasst.

Vermessen wurden ein 1er-Modul, ein 3er-Modul und ein B-Modul. Die Module
entsprechen in der vorgesehenen Konfiguration Bemessungsströmen von
\SI{250}{\ampere}, \SI{630}{\ampere} und \SI{70}{\ampere}. Für das
1er- und das 3er-Modul wurden die Prüfströme
\SI{50}{\ampere} und \SI{100}{\ampere} verwendet. Beim B-Modul betrugen
die Prüfströme \SI{10}{\ampere} und \SI{20}{\ampere}. Die Verwendung von
jeweils zwei Stromstufen ermöglicht die Prüfung, ob der Spannungsfall im
untersuchten Bereich näherungsweise proportional mit dem Strom ansteigt.

Für jeden Messpunkt wurde der Spannungsfall zunächst mit unverdrillten und
anschließend mit verdrillten Messleitungen aufgenommen. Bei Wechselstrom kann
die von den Messleitungen aufgespannte Fläche einen induzierten
Spannungsanteil erfassen. Dies ist bei den hier gemessenen Spannungsfällen im
Millivoltbereich nicht vernachlässigbar. Zur Kontrolle wurden die beiden
Messspitzen am Anschluss L1 an denselben elektrischen Punkt gehalten. Obwohl
der reale Spannungsfall zwischen den Messspitzen in diesem Fall null sein
müsste, zeigte das Messgerät abhängig von Messung und Leitungsführung Werte
zwischen \SI{1.6}{\milli\volt} und \SI{3.4}{\milli\volt}. Die Anzeige belegt
einen verbleibenden induzierten beziehungsweise störungsbedingten
Messanteil.

Der Vergleich der beiden Stromstufen zeigt den Einfluss der Leitungsführung
noch deutlicher. Bei unverdrillten Messleitungen unterscheiden sich die aus
\(\Delta U/I\) bestimmten mittleren Widerstände zwischen den beiden
Stromstufen beim 1er-Modul um \SI{7.78}{\percent}, beim 3er-Modul um
\SI{6.23}{\percent} und beim B-Modul um \SI{17.80}{\percent}. Nach dem
Verdrillen sinken diese Abweichungen auf \SI{0.22}{\percent},
\SI{2.87}{\percent} und \SI{0.20}{\percent}. Für die Modellbildung werden
daher ausschließlich die verdrillt aufgenommenen Werte der dritten Messung
verwendet.

% BILD RMCC-B -- Messablauf der dritten Messung.
% Geeignete Quellen: Mess 3.pdf, Seite 7 oder 17 (Spannungsabgriff am
% 1er-Modul), Seite 38 (verdrillte Messleitungen) sowie Seite 51 oder 79
% (Kontrollmessung mit zusammengehaltenen Messspitzen).
% \begin{figure}[H]
%     \centering
%     \begin{subfigure}[t]{0.32\textwidth}
%         \centering
%         \includegraphics[width=\textwidth]
%         {03_Ressourcen/pic/MCC_Messung3_Spannungsabgriff.jpg}
%         \caption{Spannungsabgriff}
%     \end{subfigure}
%     \hfill
%     \begin{subfigure}[t]{0.32\textwidth}
%         \centering
%         \includegraphics[width=\textwidth]
%         {03_Ressourcen/pic/MCC_Messung3_verdrillte_Leitung.jpg}
%         \caption{Verdrillte Messleitung}
%     \end{subfigure}
%     \hfill
%     \begin{subfigure}[t]{0.32\textwidth}
%         \centering
%         \includegraphics[width=\textwidth]
%         {03_Ressourcen/pic/MCC_Messung3_Kontrollmessung_L1.jpg}
%         \caption{Kontrollmessung an L1}
%     \end{subfigure}
%     \caption{Dokumentierter Messablauf der dritten Spannungsfallmessung}
%     \label{fig:kap5_mcc_messung3_ablauf}
% \end{figure}

% BILD RMCC-C -- Strompfad und Kupferverbindungen der dritten Messung.
% Geeignete Quellen: Mess 3.pdf, Seite 17 oder 29 (1er-Modul) und
% Seite 82 oder 94 (3er-Modul). Vor dem Aktivieren Stromrichtung und
% Spannungsabgriffpunkte mit dezenten Pfeilen einzeichnen.
% \begin{figure}[H]
%     \centering
%     \includegraphics[width=0.88\textwidth]
%     {03_Ressourcen/pic/MCC_Messung3_Strompfad_Messpunkte.jpg}
%     \caption{Strompfad, Kupferverbindungen und Spannungsabgriffpunkte
%     der dritten Messung}
%     \label{fig:kap5_mcc_messung3_strompfad}
% \end{figure}

Die Restanzeige bei zusammengehaltenen Messspitzen wird nicht als konstanter
Offset von den Einzelwerten abgezogen. Ein induzierter Messanteil ist von der
räumlichen Lage der Messleitungen abhängig und kann daher nicht durch eine
einzige skalare Korrektur für alle Messpunkte beschrieben werden. Die
Restanzeige wird stattdessen als Hinweis auf die verbleibende
Messunsicherheit dokumentiert. Dies ist insbesondere beim 3er-Modul zu
beachten, da dessen Spannungsfälle aufgrund des niedrigen Ersatzwiderstands
vergleichsweise klein sind.

Für jeden Phasenpfad ergibt sich der im Modell verwendete Ersatzwiderstand aus

\begin{equation}
    R_\mathrm{MCC}
    =
    \frac{\Delta U_\mathrm{MCC}}{I_\mathrm{mess}}.
    \label{eq:kap5_rmcc_bestimmung}
\end{equation}

Die Verlustleistung wird erst anschließend mit

\begin{equation}
    P_\mathrm{V,MCC}
    =
    I^2R_\mathrm{MCC}
    \label{eq:kap5_rmcc_verlustleistung}
\end{equation}

berechnet. Zur Bestimmung des Widerstands wird sie nicht gemittelt. Sie hängt
quadratisch vom Strom ab.

Zusätzlich zur phasenweisen Messung wurden die drei Phasenpfade durch die
Kupferverbindungen in Reihe geschaltet und gemeinsam gemessen. Diese
Kontrollmessung erlaubt den Vergleich des direkt gemessenen gesamten
Spannungsfalls \(U_\mathrm{3ph}\) mit der Summe der drei separat ermittelten
Phasenspannungen. Tabelle~\ref{tab:kap5_rmcc_reihenkontrolle} zeigt, dass
die Abweichung bei allen sechs Messpunkten kleiner als
\SI{0.9}{\percent} bleibt.

\begin{table}[H]
    \centering
    \caption{Kontrolle der verdrillten Messung durch Reihenschaltung der
    drei Phasenpfade}
    \label{tab:kap5_rmcc_reihenkontrolle}
    \begin{tabular}{@{}lrrrr@{}}
        \toprule
        Modul & \(I_\mathrm{mess}\) &
        \(\sum U_{L1\ldots L3}\) & \(U_\mathrm{3ph}\) & Abweichung \\
        & \si{\ampere} & \si{\milli\volt} &
        \si{\milli\volt} & \si{\percent} \\
        \midrule
        Modul 1 & \num{50}  & \num{657.4}  & \num{662.0}  & \num{0.70} \\
        Modul 1 & \num{100} & \num{1311.9} & \num{1304.0} & \num{-0.60} \\
        Modul 3 & \num{50}  & \num{97.6}   & \num{96.8}   & \num{-0.82} \\
        Modul 3 & \num{100} & \num{189.6}  & \num{188.0}  & \num{-0.84} \\
        Modul B & \num{10}  & \num{587.0}  & \num{592.0}  & \num{0.85} \\
        Modul B & \num{20}  & \num{1171.7} & \num{1174.0} & \num{0.20} \\
        \bottomrule
    \end{tabular}
\end{table}

Am 1er-Modul wurden in der dritten Messung zusätzlich Spannungsfälle über
ausgewählten Teilabschnitten des Strompfads aufgenommen. Tabelle
~\ref{tab:kap5_rmcc_teilspannungen} zeigt exemplarisch die Messung bei
\SI{100}{\ampere}. Die Summe der einzeln erfassten Bauteilspannungen ist
kleiner als der Spannungsfall des gesamten Phasenpfads, weil nicht sämtliche
internen Leiter, Anschlussstellen und Übergangskontakte als eigene
Messabschnitte aufgenommen wurden. Für das Ersatzmodell werden die
Teilspannungen deshalb nicht zu einem Widerstand addiert. Sie dienen zur
Einordnung der Entstehungsorte, während \(R_\mathrm{MCC}\) aus dem
Spannungsfall über den vollständigen Phasenpfad gebildet wird.

\begin{table}[H]
    \centering
    \caption{Spannungsfall ausgewählter Teilabschnitte des 1er-Moduls bei
    \SI{100}{\ampere} und verdrillten Messleitungen}
    \label{tab:kap5_rmcc_teilspannungen}
    \begin{tabular}{@{}lrrr@{}}
        \toprule
        Messabschnitt & \(U_{L1}\) & \(U_{L2}\) & \(U_{L3}\) \\
        & \si{\milli\volt} & \si{\milli\volt} & \si{\milli\volt} \\
        \midrule
        gesamter Phasenpfad & \num{464.4} & \num{421.1} & \num{426.4} \\
        Motorschutz-/Leistungsschalter &
        \num{169.0} & \num{164.8} & \num{155.4} \\
        Schalter & \num{70.2} & \num{79.1} & \num{68.6} \\
        SIMOCODE &
        \num{47.5} & \num{47.7} & \num{45.3} \\
        nicht einzeln erfasste Differenz &
        \num{177.7} & \num{129.5} & \num{157.1} \\
        \bottomrule
    \end{tabular}
\end{table}

Der verwendete Prüfadapter ist in der Excel-Auswertung nicht als eigener
Widerstandsanteil angesetzt. Der Modellparameter beschreibt damit den
vollständigen vermessenen Modulpfad innerhalb der festgelegten
Spannungsabgriffpunkte.

Die Excel-Auswertung bildet für jedes gemessene Modul den arithmetischen
Mittelwert der sechs mit verdrillten Leitungen bestimmten Phasenwiderstände,
also aus drei Phasen und zwei Prüfströmen. Die daraus resultierenden
Modellparameter sind in Tabelle~\ref{tab:kap5_rmcc_parameter} aufgeführt.

Die über die drei Phasen gemittelten Widerstände betragen beim 1er-Modul
\SI{4.383}{\milli\ohm} bei \SI{50}{\ampere} und
\SI{4.373}{\milli\ohm} bei \SI{100}{\ampere}. Beim 3er-Modul ergeben sich
\SI{0.651}{\milli\ohm} und \SI{0.632}{\milli\ohm}, beim B-Modul
\SI{19.567}{\milli\ohm} und \SI{19.528}{\milli\ohm}. Die Abweichungen
zwischen den beiden Stromstufen liegen damit, bezogen auf den jeweils ersten
Mittelwert, bei etwa \SI{0.22}{\percent}, \SI{2.87}{\percent} und
\SI{0.20}{\percent}. Der annähernd konstante Quotient \(\Delta U/I\)
stützt die Verwendung eines ohmschen Ersatzwerts im untersuchten
Messstrombereich.

\begin{table}[H]
    \centering
    \caption{Im Modell verwendete MCC-Ersatzwiderstände je Phase}
    \label{tab:kap5_rmcc_parameter}
    \begin{tabularx}{\textwidth}{@{}p{0.18\textwidth}
        p{0.20\textwidth}p{0.23\textwidth}X@{}}
        \toprule
        Modul & \(I_\mathrm{N}\) & \(R_\mathrm{MCC}\) je Phase & Grundlage \\
        \midrule
        Modul 1 &
        \SI{250}{\ampere} &
        \SI{4.378}{\milli\ohm} &
        Mittelwert der verdrillten Messungen bei
        \SI{50}{\ampere} und \SI{100}{\ampere} \\

        Modul 2 &
        \SI{400}{\ampere} &
        \SI{1.650}{\milli\ohm} &
        abgeschätzt aus der mittleren Bemessungsverlustleistung von
        Modul~1 und Modul~3 \\

        Modul 3 &
        \SI{630}{\ampere} &
        \SI{0.641}{\milli\ohm} &
        Mittelwert der verdrillten Messungen bei
        \SI{50}{\ampere} und \SI{100}{\ampere} \\

        Modul A &
        \SI{40}{\ampere} &
        \SI{34.208}{\milli\ohm} &
        abgeschätzt durch Skalierung des B-Moduls mit dem Verhältnis
        \(70/40\) \\

        Modul B &
        \SI{70}{\ampere} &
        \SI{19.548}{\milli\ohm} &
        Mittelwert der verdrillten Messungen bei
        \SI{10}{\ampere} und \SI{20}{\ampere} \\
        \bottomrule
    \end{tabularx}
\end{table}

Für eine phasenbezogene Modellierung können die Messwerte zusätzlich getrennt
nach Außenleitern ausgewertet werden. Dazu wird je Phase über die beiden
Prüfströme gemittelt. Tabelle~\ref{tab:kap5_rmcc_phasenwerte} zeigt diese
Werte. Die Abweichungen zwischen L1, L2 und L3 bleiben auf diese Weise
erhalten und können in einem dreiphasigen Gesamtmodell berücksichtigt werden.

\begin{table}[H]
    \centering
    \caption{Phasenbezogene Mittelwerte der gemessenen
    MCC-Ersatzwiderstände}
    \label{tab:kap5_rmcc_phasenwerte}
    \begin{tabular}{@{}lrrrr@{}}
        \toprule
        Modul & \(R_{\mathrm{MCC},L1}\) & \(R_{\mathrm{MCC},L2}\) &
        \(R_{\mathrm{MCC},L3}\) & Gesamtmittelwert \\
        & \si{\milli\ohm} & \si{\milli\ohm} & \si{\milli\ohm} &
        \si{\milli\ohm} \\
        \midrule
        Modul 1 & \num{4.611} & \num{4.255} & \num{4.268} & \num{4.378} \\
        Modul 3 & \num{0.796} & \num{0.523} & \num{0.606} & \num{0.641} \\
        Modul B & \num{19.308} & \num{18.558} & \num{20.778} & \num{19.548} \\
        \bottomrule
    \end{tabular}
\end{table}

Die Werte von Modul~1, Modul~3 und Modul~B stammen ausschließlich aus der
dritten Messung und dem zugehörigen Messprotokoll. Für Modul~2 und Modul~A
lagen keine eigenen Messreihen vor. Diese beiden Werte sind daher als
Modellannahmen und nicht als Messergebnisse zu kennzeichnen. Die Abschätzung
des 2er-Moduls basiert auf der mittleren Bemessungsverlustleistung der beiden
benachbarten großen Modulbaugrößen:

\begin{equation}
    P_{\mathrm{V,N},2}
    =
    \frac{P_{\mathrm{V,N},1}+P_{\mathrm{V,N},3}}{2},
    \qquad
    R_{\mathrm{MCC},2}
    =
    \frac{P_{\mathrm{V,N},2}}{I_{\mathrm{N},2}^{2}}.
    \label{eq:kap5_rmcc_modul2_abschaetzung}
\end{equation}

Für das A-Modul wurde

\begin{equation}
    R_{\mathrm{MCC,A}}
    =
    R_{\mathrm{MCC,B}}
    \frac{I_{\mathrm{N,B}}}{I_{\mathrm{N,A}}}
    \label{eq:kap5_rmcc_modula_abschaetzung}
\end{equation}

angesetzt. Die Gleichung stellt eine Skalierungsannahme und kein allgemeines
Gesetz für MCC-Module dar.

Für das im thermischen Modell betrachtete \SI{630}{\ampere}-Modul wird der
gemessene Wert

\begin{equation}
    R_{\mathrm{MCC},630}
    =
    \SI{0.641}{\milli\ohm}
    \label{eq:kap5_rmcc_630}
\end{equation}

verwendet. Im reduzierten Modell wird der aus dem Spannungsfall berechnete
Wert als ohmscher Ersatzwiderstand behandelt. Da keine Phasenwinkelmessung
vorliegt, ist diese Reduktion als Modellannahme zu verstehen. Das Verdrillen
der Messleitungen reduziert einen möglichen induzierten Messanteil, ersetzt
jedoch keine getrennte Bestimmung von Wirk- und Blindanteil. Für den zunächst
repräsentativ betrachteten einphasigen 630-A-Lastzweig wird der
Gesamtmittelwert verwendet. Wird das Gesamtmodell phasenbezogen ausgeführt,
sind stattdessen die in Tabelle~\ref{tab:kap5_rmcc_phasenwerte} angegebenen
Einzelwerte einzusetzen.

Die Widerstände wurden bei \SI{50}{\ampere} und \SI{100}{\ampere}
bestimmt. Ihre Übertragung auf den Bemessungsstrom setzt voraus, dass der
Modulpfad im betrachteten Bereich näherungsweise linear und sein Widerstand
während der Messung annähernd konstant bleibt. Eine mögliche Eigenerwärmung des
MCC-Moduls bei Bemessungsstrom wird durch den konstanten Ersatzwiderstand
nicht beschrieben. Das thermische Teilmodell dieses Kapitels bezieht sich
ausschließlich auf den V2A-Widerstand.


% -----------------------------------------------------------------------------
% 5.4 Thermisches Grundmodell -- Kreis 1
% -----------------------------------------------------------------------------

\subsection{Thermisches Grundmodell -- Kreis 1}
\label{sec:thermisches_modell_kreis1}

Nach der Zerlegung wurde das thermische Verhalten zunächst unabhängig vom
elektrischen Lastzweig aufgebaut. Kreis~1 bildet einen thermischen Speicher
erster Ordnung. Die zugeführte elektrische Verlustleistung
\(P_\mathrm{el}\) wird als Wärmestrom interpretiert. Ein Teil dieser Leistung
erhöht die im V2A-Körper gespeicherte Wärmeenergie, während der verbleibende
Anteil über den thermischen Widerstand an die Umgebung abgegeben wird.

Die Energiebilanz lautet

\begin{equation}
    C_\mathrm{th}\frac{\mathrm{d}\Delta\vartheta}{\mathrm{d}t}
    =
    P_\mathrm{el}
    -
    \frac{\Delta\vartheta}{R_\mathrm{th}}.
    \label{eq:kap5_kreis1_dgl}
\end{equation}

Kreis~1 wurde zunächst normiert mit
\(C_\mathrm{th}=1\), \(R_\mathrm{th}=1\) und
\(P_\mathrm{el}=1\) untersucht. Diese normierte Parametrierung
diente nicht der Beschreibung des realen V2A-Widerstands. Mit ihr konnten
Vorzeichen, Rückführungsrichtung, Integrator und stationärer Grenzwert ohne
zusätzliche Parameterabhängigkeiten kontrolliert werden. Für den normierten
Fall ergeben sich

\begin{equation}
    \tau_\mathrm{th}=R_\mathrm{th}C_\mathrm{th}=1
    \label{eq:kap5_kreis1_tau}
\end{equation}

und

\begin{equation}
    \Delta\vartheta_\infty
    =
    P_\mathrm{el}R_\mathrm{th}
    =1.
    \label{eq:kap5_kreis1_endwert}
\end{equation}

Neben der Temperatur wurden die gespeicherte beziehungsweise zur
Temperaturänderung wirksame Leistung und die abgeführte Wärmeleistung
ausgegeben. Im stationären Zustand muss die zeitliche Temperaturänderung
verschwinden. Damit gilt

\begin{equation}
    \dot Q_\mathrm{C}=0,
    \qquad
    P_\mathrm{el}=\dot Q_\mathrm{ab}.
    \label{eq:kap5_kreis1_stationaer}
\end{equation}

% BILD 5.7 -- liegt als RC_Modell_1.pdf vor.
% \begin{figure}[H]
%     \centering
%     \includegraphics[page=3,width=\textwidth,
%         trim={3.2cm 4.0cm 3.2cm 4.0cm},clip]
%     {03_Ressourcen/pic/RC_Modell_1.pdf}
%     \caption{Kreis~1: normiertes thermisches RC-Grundmodell
%     (eigene Darstellung in MATLAB/Simulink)}
%     \label{fig:kap5_rc_modell_1}
% \end{figure}

Die erfolgreiche Prüfung dieses normierten Modells war die Voraussetzung für
die Übernahme der Struktur in den physikalisch parametrierten Kreis~2.


% -----------------------------------------------------------------------------
% 5.5 Physikalische Parametrierung -- Kreis 2
% -----------------------------------------------------------------------------

\subsection{Physikalische Parametrierung -- Kreis 2}
\label{sec:thermisches_modell_kreis2}

In Kreis~2 wurde die normierte Struktur durch die geometrischen und
thermischen Parameter des V2A-Widerstands ersetzt. Die elektrische
Verlustleistung blieb zunächst eine unabhängig vorgegebene Eingangsgröße.
Dadurch konnte die thermische Strecke geprüft werden, ohne dass eine
gleichzeitig wirkende elektrische Rückkopplung das Ergebnis verändert.

Die Masse des V2A-Körpers wird aus Geometrie und Dichte bestimmt:

\begin{equation}
    m
    =
    \rho_\mathrm{m}\,l\,b\,d.
    \label{eq:kap5_masse_v2a}
\end{equation}

Unter Vernachlässigung der beiden kleinen Stirnflächen wird für die wirksame
Oberfläche der Längs- und Schmalseiten

\begin{equation}
    A_\mathrm{O}
    =
    2l(b+d)
    \label{eq:kap5_oberflaeche_v2a}
\end{equation}

angesetzt. Die thermische Kapazität folgt aus Masse und spezifischer
Wärmekapazität:

\begin{equation}
    C_\mathrm{th}
    =
    m c_\mathrm{p}
    =
    \rho_\mathrm{m}\,l\,b\,d\,c_\mathrm{p}.
    \label{eq:kap5_cth}
\end{equation}

Aus der wirksamen Oberfläche \(A_\mathrm{O}\) und dem angesetzten
Wärmeübergangskoeffizienten \(\alpha_\mathrm{K}\) folgt der thermische
Widerstand

\begin{equation}
    R_\mathrm{th}
    =
    \frac{1}{\alpha_\mathrm{K}A_\mathrm{O}}
    \label{eq:kap5_rth}
\end{equation}

angenähert. Für den untersuchten \SI{630}{\ampere}-Lastzweig werden
\(l=\SI{4.0}{\metre}\), \(b=\SI{150}{\milli\metre}\) und
\(d=\SI{3}{\milli\metre}\) angesetzt. Mit
\(\rho_\mathrm{m}=\SI{7900}{\kilogram\per\cubic\metre}\),
\(c_\mathrm{p}=\SI{500}{\joule\per\kilogram\per\kelvin}\) und
\(\alpha_\mathrm{K}=\SI{7.5}{\watt\per\square\metre\per\kelvin}\)
ergeben sich die in Tabelle~\ref{tab:kap5_thermische_parameter_630}
zusammengefassten Werte.

\begin{table}[H]
    \centering
    \caption{Parameter des thermischen Modells für den untersuchten
    \SI{630}{\ampere}-V2A-Widerstand}
    \label{tab:kap5_thermische_parameter_630}
    \begin{tabular}{@{}lll@{}}
        \toprule
        Größe & Formelzeichen & Wert \\
        \midrule
        elektrischer Kaltwiderstand &
        \(R_{\mathrm{V2A},20}\) &
        \SI{6.489}{\milli\ohm} \\

        Masse des V2A-Körpers &
        \(m\) &
        \SI{14.22}{\kilogram} \\

        wirksame Oberfläche &
        \(A_\mathrm{O}\) &
        \SI{1.224}{\square\metre} \\

        thermische Kapazität &
        \(C_\mathrm{th}\) &
        \SI{7110}{\joule\per\kelvin} \\

        thermischer Widerstand &
        \(R_\mathrm{th}\) &
        \SI{0.1089}{\kelvin\per\watt} \\

        thermische Zeitkonstante &
        \(\tau_\mathrm{th}\) &
        \SI{774.5}{\second} \\

        fünffache Zeitkonstante &
        \(5\tau_\mathrm{th}\) &
        \SI{3872.6}{\second} \\

        Umgebungstemperatur &
        \(\vartheta_\mathrm{u}\) &
        \SI{20}{\degreeCelsius} \\
        \bottomrule
    \end{tabular}
\end{table}

Die Temperaturantwort auf einen konstanten Leistungssprung
\(P_0\) lautet für das lineare RC-Modell

\begin{equation}
    \Delta\vartheta(t)
    =
    P_0R_\mathrm{th}
    \left(1-\mathrm{e}^{-t/\tau_\mathrm{th}}\right).
    \label{eq:kap5_thermische_sprungantwort}
\end{equation}

Nach einer Zeitkonstante sind etwa \(63.2\,\%\), nach fünf
Zeitkonstanten etwa \(99.3\,\%\) des stationären Endwerts erreicht. Diese
Beziehungen ermöglichen eine unabhängige Kontrolle der numerischen
Simulation. Der stationäre Endwert

\begin{equation}
    \Delta\vartheta_\infty
    =
    P_0R_\mathrm{th}
    \label{eq:kap5_thermischer_endwert}
\end{equation}

gilt unmittelbar nur bei konstanter Eingangsleistung. Im später gekoppelten
Modell verändert sich die Verlustleistung mit Strom und Widerstand.

Der Leistungssprung ist ein Sonderfall der in
Abschnitt~\ref{sec:thermisches_rc_grundlagen} hergeleiteten Faltung. Für
Kreis~2 mit konstanten thermischen Parametern und einer von außen
vorgegebenen Verlustleistung kann jeder beliebige Leistungsverlauf mit der
Gewichtsfunktion des PT1-Glieds auf die Übertemperatur abgebildet werden.
Jeder zurückliegende Leistungsanteil wirkt dabei exponentiell abklingend auf
die aktuelle Temperatur. In Kreis~3 bleibt diese Beschreibung für das
isolierte thermische Teilmodell gültig. Durch die Rückkopplung von
Temperatur, Widerstand und Verlustleistung ist das elektrothermische
Gesamtsystem jedoch nichtlinear und kann nicht mehr allein durch eine Faltung
mit einem vorab bekannten Leistungsverlauf gelöst werden.

% BILD 5.8 -- liegt als RC_Modell_2.pdf vor.
% \begin{figure}[H]
%     \centering
%     \includegraphics[width=\textwidth,
%         trim={2.5cm 1.5cm 1.0cm 0.5cm},clip]
%     {03_Ressourcen/pic/RC_Modell_2.pdf}
%     \caption{Kreis~2: thermisches Teilmodell mit physikalischen Parametern
%     und vorgegebener Verlustleistung
%     (eigene Darstellung in MATLAB/Simulink)}
%     \label{fig:kap5_rc_modell_2}
% \end{figure}

Der Wärmeübergangskoeffizient wurde nicht durch eine zeitaufgelöste
Temperaturmessung identifiziert. Die berechnete Temperatur ist daher eine
Modellgröße des gewählten Parametersatzes. Kreis~2 dient vorrangig zur
Plausibilisierung der Struktur, der Zeitkonstante und der Energiebilanz.


% -----------------------------------------------------------------------------
% 5.6 Elektrothermische Rückkopplung -- Kreis 3
% -----------------------------------------------------------------------------

\subsection{Elektrothermische Rückkopplung -- Kreis 3}
\label{sec:thermisches_modell_kreis3}

Kreis~3 schließt den elektrothermischen Wirkzusammenhang. Der elektrische
Widerstand des V2A-Materials wird als linear temperaturabhängig beschrieben:

\begin{equation}
    R_\mathrm{V2A}(\vartheta)
    =
    R_{\mathrm{V2A},20}
    \left[
        1+\alpha_\mathrm{el}
        \left(\vartheta-\SI{20}{\degreeCelsius}\right)
    \right].
    \label{eq:kap5_rv2a_temperatur}
\end{equation}

Der Kaltwiderstand folgt mit
\(\rho_{\mathrm{el},20}=0.73\,\Omega\,\mathrm{mm^2/m}\) aus

\begin{equation}
    R_{\mathrm{V2A},20}
    =
    \rho_{\mathrm{el},20}\frac{l}{bd}
    =
    \SI{6.489}{\milli\ohm}.
    \label{eq:kap5_rv2a_20}
\end{equation}

In dieser Modellstufe lag die vorgegebene Spannung unmittelbar am
V2A-Widerstand. Strom und Verlustleistung wurden deshalb mit

\begin{align}
    I_\mathrm{V2A}(t)
    &=
    \frac{U_\mathrm{V2A}(t)}{R_\mathrm{V2A}(\vartheta(t))},
    \\
    P_\mathrm{V2A}(t)
    &=
    I_\mathrm{V2A}^{2}(t)R_\mathrm{V2A}(\vartheta(t))
    =
    \frac{U_\mathrm{V2A}^{2}(t)}{R_\mathrm{V2A}(\vartheta(t))}
    \label{eq:kap5_kreis3_elektrisch}
\end{align}

bestimmt. Die Verlustleistung wird dem thermischen Modell zugeführt. Die
daraus berechnete Temperatur verändert wiederum den elektrischen Widerstand
und damit Strom und Leistung. Es entsteht die Wirkungskette

\begin{equation}
    U_\mathrm{V2A}
    \longrightarrow
    I_\mathrm{V2A}
    \longrightarrow
    P_\mathrm{V2A}
    \longrightarrow
    \vartheta
    \longrightarrow
    R_\mathrm{V2A}(\vartheta)
    \longrightarrow
    I_\mathrm{V2A}.
    \label{eq:kap5_kreis3_wirkungskette}
\end{equation}

Für diese Untersuchung wurde eine Spannungssprungfunktion verwendet. Der
Sprung besitzt gegenüber der realitätsnäheren Rampe einen methodischen
Vorteil: Der Beginn der elektrischen Belastung ist eindeutig festgelegt und
die daraus folgende thermische Antwort kann ohne überlagerten
Spannungshochlauf beurteilt werden. Außerdem lässt sich das entkoppelte
thermische Grundmodell unmittelbar mit der bekannten PT1-Sprungantwort
vergleichen. Der Spannungssprung stellt daher eine Prüfanregung des
thermischen Modells und keine Nachbildung des realen
Säulenstelltransformators dar.

% OFFEN:
% Sprungzeitpunkt, U_sprung, Solver und Schrittweite erst aus dem final
% verwendeten Kreis-3-Modell übernehmen. Zwischenstände nicht vermischen.

% BILD 5.9 -- Gesamtansicht, Seite 1 der Datei RC_Modell_3.pdf.
% \begin{figure}[H]
%     \centering
%     \includegraphics[page=1,width=\textwidth]
%     {03_Ressourcen/pic/RC_Modell_3.pdf}
%     \caption{Kreis~3: elektrothermisch gekoppeltes Modell mit
%     Spannungssprung und temperaturabhängigem V2A-Widerstand
%     (eigene Darstellung in MATLAB/Simulink)}
%     \label{fig:kap5_rc_modell_3}
% \end{figure}

% BILD 5.10 -- thermisches Subsystem, Seite 3 der Datei RC_Modell_3.pdf.
% \begin{figure}[H]
%     \centering
%     \includegraphics[page=3,width=\textwidth]
%     {03_Ressourcen/pic/RC_Modell_3.pdf}
%     \caption{Thermisches Subsystem des elektrothermisch gekoppelten
%     Kreis-3-Modells (eigene Darstellung in MATLAB/Simulink)}
%     \label{fig:kap5_rc_modell_3_subsystem}
% \end{figure}

Kreis~3 ist noch nicht mit dem zusammengeführten Gesamtmodell
gleichzusetzen. Die Spannungsquelle bildet in diesem Stand nicht den linearen
Hochlauf mit OSG ab. Außerdem fehlt der gemessene Ersatzwiderstand des
MCC-Moduls. Die Ergebnisse von Kreis~3 beschreiben daher ausschließlich den
elektrothermischen V2A-Lastwiderstand unter der verwendeten Sprunganregung.
Sie dürfen nicht als vollständige Vorhersage des realen Prüfstands
interpretiert werden.


% -----------------------------------------------------------------------------
% 5.7 Zusammenführung der Teilmodelle
% -----------------------------------------------------------------------------

\subsection{Zusammenführung der Teilmodelle}
\label{sec:zusammenfuehrung_gesamtmodell}

Das zusammengeführte Gesamtmodell entsteht aus den zuvor getrennt
untersuchten und plausibilisierten Teilmodellen. Es dient als Referenzmodell
des bestehenden Lastzweigs, erhebt jedoch nicht den Anspruch, sämtliche
Eigenschaften des realen Prüfstands abzubilden. Für einen betrachteten
Außenleiter des \SI{630}{\ampere}-Abgangs werden dazu

\begin{itemize}
    \item das Quellenmodell nach Abschnitt~\ref{sec:spannungsquelle_osg},
    \item der gemessene MCC-Ersatzwiderstand
          \(R_\mathrm{MCC,630}=\SI{0.641}{\milli\ohm}\),
    \item der geometrisch bestimmte V2A-Kaltwiderstand
          \(R_\mathrm{V2A,20}=\SI{6.489}{\milli\ohm}\) und
    \item das in Kreis~1 und Kreis~2 geprüfte thermische RC-Modell
\end{itemize}

verwendet. Der Strom des Zweigs wird im ohmsch reduzierten Modell durch

\begin{equation}
    I(t)
    =
    \frac{U_\mathrm{eff,Zeiger}(t)}
    {R_\mathrm{MCC,630}+R_\mathrm{V2A}(\vartheta(t))}
    \label{eq:kap5_gesamtmodell_strom}
\end{equation}

bestimmt. Weitere Leitungs- und Kontaktwiderstände dürfen nur dann zusätzlich
angesetzt werden, wenn sie nicht bereits in der Spannungsfallmessung des
Moduls enthalten sind.

Für das thermische Modell ist ausschließlich die im V2A-Widerstand
umgesetzte Leistung maßgebend:

\begin{equation}
    P_\mathrm{V2A}(t)
    =
    I^2(t)R_\mathrm{V2A}(\vartheta(t)).
    \label{eq:kap5_gesamtmodell_pv2a}
\end{equation}

Die Verlustleistung
\(I^2R_\mathrm{MCC}\) entsteht innerhalb der MCC-Abgangseinheit und liegt
außerhalb der thermischen Systemgrenze des V2A-Schienenmodells. Sie wird daher
nicht auf die thermische Kapazität des V2A-Widerstands aufgeschaltet.

Die elektrische Beschreibung und die thermische Energiebilanz ergeben
gemeinsam das folgende gekoppelte Gleichungssystem:

\begin{align}
    I(t)
    &=
    \frac{U_\mathrm{eff,Zeiger}(t)}
    {R_\mathrm{MCC,630}+R_\mathrm{V2A}(\vartheta(t))},
    \\
    P_\mathrm{V2A}(t)
    &=
    I^2(t)R_\mathrm{V2A}(\vartheta(t)),
    \\
    C_\mathrm{th}\frac{\mathrm{d}\Delta\vartheta(t)}{\mathrm{d}t}
    &=
    P_\mathrm{V2A}(t)
    -
    \frac{\Delta\vartheta(t)}{R_\mathrm{th}},
    \\
    R_\mathrm{V2A}(\vartheta(t))
    &=
    R_\mathrm{V2A,20}
    \left[1+\alpha_\mathrm{el}\Delta\vartheta(t)\right].
    \label{eq:kap5_gesamtmodell_gekoppelt}
\end{align}

Eine wichtige Abgrenzung betrifft den Bezugspunkt der Spannung. In Kreis~3
wurde \(U_\mathrm{V2A}\) unmittelbar über dem V2A-Widerstand vorgegeben. Bei
\(R_\mathrm{V2A,20}=\SI{6.489}{\milli\ohm}\) führen etwa
\SI{4.089}{\volt} rechnerisch zu \SI{630}{\ampere}. Wird der gemessene
MCC-Ersatzwiderstand in Reihe ergänzt, liegt dieselbe Spannung nun vor der
Reihenschaltung. Der kalte Strom beträgt dann nur

\begin{equation}
    I(\SI{20}{\degreeCelsius})
    =
    \frac{\SI{4.089}{\volt}}
    {\SI{0.641}{\milli\ohm}+\SI{6.489}{\milli\ohm}}
    \approx
    \SI{573.5}{\ampere}.
    \label{eq:kap5_gesamtmodell_strom_4089v}
\end{equation}

Soll der kalte Zweig unter denselben Vereinfachungen
\SI{630}{\ampere} führen, muss die vor der Reihenschaltung erforderliche
Effektivspannung etwa

\begin{equation}
    U_\mathrm{eff,erf}
    =
    \SI{630}{\ampere}
    \left(
        \SI{0.641}{\milli\ohm}
        +
        \SI{6.489}{\milli\ohm}
    \right)
    \approx
    \SI{4.49}{\volt}
    \label{eq:kap5_gesamtmodell_spannung_erforderlich}
\end{equation}

betragen. Vor der endgültigen Simulation ist deshalb festzulegen, an welcher
Stelle des realen und simulierten Zweigs die angegebene Spannung erfasst wird.
Erst danach kann der Endwert der linearen Quellenrampe eindeutig festgelegt
werden.

% BILD 5.11 -- erst nach der tatsächlichen Zusammenführung exportieren.
% \begin{figure}[H]
%     \centering
%     \includegraphics[width=\textwidth]
%     {03_Ressourcen/pic/Gesamtmodell_Referenzlastzweig_630A.pdf}
%     \caption{Zusammengeführtes Modell des betrachteten
%     \SI{630}{\ampere}-Lastzweigs mit linearer OSG-Spannungsquelle,
%     gemessenem MCC-Ersatzwiderstand und elektrothermischer Rückkopplung
%     (eigene Darstellung in MATLAB/Simulink)}
%     \label{fig:kap5_gesamtmodell_referenzlastzweig}
% \end{figure}

% OFFEN -- nach Fertigstellung des Gesamtmodells ergänzen:
% - tatsächlicher Spannungsendwert und dessen Messort,
% - Simulationszeit, Solver und Schrittweite,
% - Anfangs- und Endstrom,
% - Temperatur-, Widerstands- und Leistungsverlauf,
% - Vergleich mit Kreis 3,
% - Ableitung der erforderlichen Stellreserve.


% -----------------------------------------------------------------------------
% 5.8 Plausibilisierung des zusammengefuehrten Gesamtmodells
%     und Aussagegrenzen
% -----------------------------------------------------------------------------

\subsection{Plausibilisierung des zusammengeführten Gesamtmodells und
Aussagegrenzen}
\label{sec:kap5_plausibilisierung_grenzen}

Die Plausibilisierung erfolgt auf zwei Ebenen. Zunächst werden die Teilmodelle
jeweils mit Beziehungen geprüft, die unabhängig von der numerischen Umsetzung
gelten. Nach ihrer Zusammenführung wird zusätzlich geprüft, ob die
Schnittstellen und die geschlossene elektrothermische Wirkungskette des
Gesamtmodells konsistent sind. Für die Spannungsquelle muss der Scheitelwert
im stationären Zustand

\begin{equation}
    \hat U=\sqrt{2}\,U_\mathrm{eff}
    \label{eq:kap5_plausi_scheitelwert}
\end{equation}

erfüllen. Der Effektivwertzeiger muss der vorgegebenen Rampe folgen, während
die Abweichung der klassischen RMS-Auswertung durch deren Messfenster
erklärbar sein muss. Für die MCC-Messwerte ist zu prüfen, ob sich bei beiden
Prüfströmen vergleichbare Werte von \(R_\mathrm{MCC}=\Delta U/I\) ergeben.
Damit wird die angenommene Proportionalität des Spannungsfalls im untersuchten
Strombereich kontrolliert.

Für Kreis~1 und Kreis~2 werden stationärer Endwert und Zeitkonstante mit den
analytischen Beziehungen verglichen. Ohne zugeführte Verlustleistung und bei
einer Anfangstemperatur gleich der Umgebungstemperatur darf keine
Temperaturänderung auftreten. Bei konstanter Eingangsleistung müssen nach
\(\tau_\mathrm{th}\) etwa \(63.2\,\%\) und nach
\(5\tau_\mathrm{th}\) etwa \(99.3\,\%\) des Endwerts erreicht sein.

Für Kreis~3 und das zusammengeführte Gesamtmodell muss die Wirkungskette
konsistent sein: Ein positiver Temperaturkoeffizient führt bei konstanter
Spannung zu einem Anstieg des V2A-Widerstands und damit zu einer Abnahme des
Stroms. Wird \(\alpha_\mathrm{el}=0\) gesetzt, muss dieser thermisch
verursachte Stromdrift verschwinden. Die thermische Eingangsleistung muss
stets aus dem V2A-Anteil und nicht aus der gesamten Serienverlustleistung
gebildet werden. Im zusammengeführten Gesamtmodell ist außerdem zu prüfen,
ob der kalte Anfangsstrom mit der vorgegebenen Quellenspannung und der Summe
aus \(R_\mathrm{MCC}\) und \(R_\mathrm{V2A,20}\) übereinstimmt.

Die Herkunft der maßgebenden Parameter ist in
Tabelle~\ref{tab:kap5_parameterherkunft} zusammengefasst.

\begin{table}[H]
    \centering
    \caption{Herkunft und Einordnung der maßgebenden Modellparameter}
    \label{tab:kap5_parameterherkunft}
    \begin{tabularx}{\textwidth}{@{}p{0.26\textwidth}
        p{0.29\textwidth}X@{}}
        \toprule
        Parameter & Herkunft & Einordnung \\
        \midrule
        \(U_\mathrm{eff}(t)\), \(f\) &
        definierter Quellenfall und Prüfstandsfrequenz &
        Modellanregung; Spannungsendwert im Gesamtmodell noch abzugleichen \\

        \(R_\mathrm{MCC,630}\) &
        dritte Spannungsfallmessung mit Fotodokumentation und Messprotokoll &
        messgestützter Ersatzwiderstand \\

        \(R_\mathrm{MCC,400}\), \(R_\mathrm{MCC,A}\) &
        Excel-basierte Abschätzungen &
        Modellannahmen, keine Messergebnisse \\

        \(R_\mathrm{V2A,20}\) &
        Geometrie und spezifischer elektrischer Widerstand &
        rechnerisch bestimmter Kaltwiderstand \\

        \(\alpha_\mathrm{el}\), \(\rho_\mathrm{m}\), \(c_\mathrm{p}\) &
        Werkstoffkennwerte &
        auf das reduzierte Modell übertragene Kennwerte \\

        \(R_\mathrm{th}\) &
        Oberfläche und angenommener Wärmeübergangskoeffizient &
        Modellparameter mit hoher Unsicherheit \\

        \(C_\mathrm{th}\) &
        Geometrie, Dichte und spezifische Wärmekapazität &
        rechnerisch bestimmter Modellparameter \\
        \bottomrule
    \end{tabularx}
\end{table}

Das thermische Modell verwendet eine räumlich homogene Ersatztemperatur.
Lokale Heißstellen, Wärmeleitung über Anschlüsse und Halterungen,
temperaturabhängige Konvektion, Wärmestrahlung und die gegenseitige Erwärmung
benachbarter Schienen werden nicht separat abgebildet. Insbesondere der
angesetzte Wärmeübergangskoeffizient wurde nicht anhand einer vollständigen
Temperaturmessreihe identifiziert. Berechnete Endtemperaturen sind deshalb als
Ergebnisse des jeweiligen Rechenfalls und nicht als genaue Vorhersage des
realen Prüfstands zu interpretieren.

Auch die Spannungsfallmessung bestimmt ohne zusätzliche Phasenmessung keinen
getrennten Wirk- und Blindanteil. Die Verwendung als ohmscher
Ersatzwiderstand ist für das reduzierte Modell begründet zu dokumentieren.
Eine Plausibilisierung anhand analytischer Grenzfälle weist die innere
Konsistenz der Modellgleichungen nach, ersetzt jedoch keine experimentelle
Validierung des vollständig gekoppelten Gesamtmodells.

Die elektrische Parametrierung, die thermische Rückkopplung und ihre
Schnittstellen sind damit festgelegt. Die Endwerte aus Kreis~3 unter
Sprunganregung sind jedoch keine Ergebnisse des zusammengeführten
Gesamtmodells. Dessen quantitative Plausibilisierung folgt nach Einbau der
linearen Spannungsquelle, Festlegung des Spannungsbezugspunkts und Übernahme
des gemessenen MCC-Ersatzwiderstands. Sie umfasst den Vergleich von
Anfangsstrom, Temperatur, Widerstand, Leistung und thermischer Energiebilanz
mit den zuvor hergeleiteten Beziehungen.

% ABSCHLIESSENDER ERGEBNISTEIL:
% Nach der finalen Simulation hier nicht einfach weitere Theorie ergänzen.
% Stattdessen zwei bis vier Ergebnisabbildungen einfügen und jeweils angeben:
% 1. verwendete Randbedingungen,
% 2. dargestellte Größen,
% 3. Anfangs- und Endwerte,
% 4. physikalische Ursache des Verlaufs,
% 5. Folgerung für die erforderliche Nachführung.
